\section{Introduction}
\label{sec:introduction}

A model that predicts average human behavior may still fit some populations better than others. Centaur adapts a language model to behavioral experiments and generalizes across participants and tasks \citep{binz2025centaur}. We ask how the population represented in additional fine-tuning data shapes its predictions.

Fine-tuning may change predictions without improving them. It may learn a general response preference, or how a population responds to particular items. Either change may fail to transfer faithfully to another task. These outcomes require separate tests.

We use US and Chinese samples from Giving and Taking dictator games \citep{wang2026decomposing}, category-versus-relation word choices \citep{cao2024differences}, and explicit lotteries \citep{anllo2024preferences}. Within each domain, US-only, China-only, and pooled fine-tuning use equal data budgets and identical prompts. Evaluation holds out people and complete items. Giving also provides a transfer test: adapters fine-tuned only on Giving are evaluated on Taking.

We ask four questions: Does frozen Centaur fit one population better? Does fine-tuning population change unseen-item predictions? Do these changes improve fit and transfer faithfully? Does fine-tuning outperform a simple response adjustment? Our contribution is to separate these outcomes empirically. Country comparisons do not identify a cultural mechanism.
